% ============================================================================
% sec5_empirical_strategy.tex — Paper 2: Estimation strategy
% ============================================================================

\section{Estimation Strategy}
\label{sec:empirical}

The estimation has four components, ordered by their structural content.
\emph{First}, we estimate the mixture model from
Section~\ref{sec:likelihood}, which identifies the cover-bid
distribution and selects between regimes
(Section~\ref{sec:structural_estimation}). \emph{Second}, we test
the model's predictions using reduced-form regressions, matching
estimators, and an IV bracketing exercise
(Section~\ref{sec:reduced_form}). \emph{Third}, we validate detection
performance against CADE convictions and compare with bid-level
screens (Section~\ref{sec:detection_validation}). \emph{Fourth}, we
implement Bajari--Ye bid-coordination tests using the FL partition
(Section~\ref{sec:bajari_ye}).

% ── 5.1 Structural Estimation ────────────────────────────────────

\subsection{Structural Likelihood Estimation}
\label{sec:structural_estimation}

The estimation proceeds in two stages on the bid-level sample
(23.2~million genuine losing bids; 189,381 FL losing bids).

\paragraph{Stage 1: Genuine bid distribution.}
We estimate $(\boldsymbol{\beta}, \sigma_g^2)$ from non-FL losing bids
by maximum likelihood under the log-normal
specification~\eqref{eq:genuine_dist}:
\begin{equation}
  (\hat{\boldsymbol{\beta}}, \hat{\sigma}_g^2) =
  \operatorname*{arg\,max}_{\boldsymbol{\beta}, \sigma_g^2}
  \sum_{t=1}^{T} \sum_{j \in \mathcal{G}_t}
  \log \phi\!\left(
    \frac{\log b_{jt} - \mathbf{x}_j'\boldsymbol{\beta}
    - \hat{\alpha}_i - \hat{\lambda}_t}{\sigma_g}
  \right)
  \label{eq:mle_genuine}
\end{equation}
In practice, this reduces to OLS of $\log b_{jt}$ on $\mathbf{x}_j$
with item and year fixed effects, yielding $\hat{\sigma}_g^2$ as the
residual variance.

\paragraph{Stage 2: Cover bid distribution.}
Conditional on $b^*_t$, we estimate cover-bid parameters from FL bids.
Under Regime~1: $\hat{\delta} = \max_{j \in \mathcal{F}_t, t}
(b_{jt} - b^*_t)$. Under Regime~2:
\begin{equation}
  \hat{\epsilon} = \overline{\log b_{jt} - \log b^*_t}, \qquad
  \hat{\sigma}_c^2 = \frac{1}{|\mathcal{F}|}
  \sum_{t,j \in \mathcal{F}_t}
  \left(\log b_{jt} - \log b^*_t - \hat{\epsilon}\right)^2
  \label{eq:mle_cover}
\end{equation}

\paragraph{Model selection.}
BIC and Kolmogorov--Smirnov goodness-of-fit jointly select the
regime:
\begin{equation}
  \text{BIC}_r = -2\ell_r + k_r \log N, \quad r \in \{1, 2\}
  \label{eq:bic}
\end{equation}

\paragraph{Price-association consistency check.}
The $n$-conditional price association compares FL-present and
FL-absent tenders, conditioning on the number of genuine competitors
$n$:
\begin{equation}
  \hat{\mu} = E[\log p \mid \text{FL present}, n] -
  E[\log p \mid \text{FL absent}, n]
  \label{eq:structural_markup}
\end{equation}
This is a calibration check: if the structural model is correctly
specified, the $n$-conditional coefficient should be consistent with
the unconditional OLS estimate, because the model predicts that
differential genuine competition does not drive the FL--price
association. The structural parameters
$(\hat{\epsilon}, \hat{\sigma}_c)$ characterize the cover-bid
\emph{mechanism}; the price association comes from the reduced-form
comparison, disciplined by the model's predictions.

\paragraph{Standard errors.}
Parametric bootstrap with 500 replications, resampling tenders to
preserve within-tender correlation.

% ── 5.2 Calibration of Cartel Primitives ──────────────────────────

\subsection{Calibration of Cartel Primitives}
\label{sec:calibration}

The mixture model identifies the cover-bid distribution
($\hat{\epsilon}$, $\hat{\sigma}_c$) but not the cartel's
decision parameters ($c_1$, $\phi_0$) that govern cover-bidder
deployment. We recover these from the model's comparative statics
using a minimum-distance calibration.

\paragraph{Moment conditions.}
From Proposition~\ref{prop:optimal_m}, the interior-solution FOC
implies:
\begin{equation}
  m^*_{kt} = \frac{\pi'(n_{kt})}{c_1 + \theta_k \phi_0}
  \label{eq:foc_m}
\end{equation}
where $\pi'(n) \equiv \partial \pi / \partial m$ is the marginal
return to cover bidding (A2 ensures concavity in $m$; the sign
of $\partial \pi' / \partial n$ is an empirical question---see
Proposition~\ref{prop:market_selection}). Taking logs:
\begin{equation}
  \log m^*_{kt} = \log \pi'(n_{kt}) - \log(c_1 + \theta_k \phi_0)
  \label{eq:log_foc}
\end{equation}

We parameterize $\pi'(n) = \pi_0 \cdot n^{\gamma}$ and
$\theta_k = \bar{\theta} \cdot \text{PBUsize}_k^{\psi}$ (detection
probability increasing in oversight capacity). The sign of $\gamma$
is unrestricted: $\gamma > 0$ implies strategic complementarity
between cover bidding and genuine competition
(Proposition~\ref{prop:market_selection}), while $\gamma < 0$ implies
substitution. This yields a log-linear specification:
\begin{equation}
  \log \bar{m}_k = a_0 + \gamma \log \bar{n}_k
  - \log\!\left(c_1 + \bar{\theta} \cdot \text{PBUsize}_k^{\psi}
  \cdot \phi_0\right) + \nu_k
  \label{eq:calibration_moment}
\end{equation}
where $\bar{m}_k$ and $\bar{n}_k$ are PBU-level means of FL count
and genuine-bidder count, respectively, and $\nu_k$ captures
unobserved PBU heterogeneity.

\paragraph{Estimation.}
We estimate Eq.~\eqref{eq:calibration_moment} on PBU-level moments
from 913 PBUs with $\geq$10 interior-solution tenders ($m^* >
\underline{n}(\tau) - n$). Because the five-parameter NLS surface
exhibits near-collinearity between $c_1$ and $\phi_0$ (enforcement
variation across PBUs is limited), we adopt a profile-likelihood
approach: $\psi$ is concentrated out over a grid
$\{0.1, 0.2, \ldots, 3.0\}$, and the remaining four parameters
$(\pi_0, \gamma, c_1, \phi_0)$ are estimated by weighted NLS at each
grid point. The profile-best $\psi^*$ minimizes the weighted residual
sum of squares. Standard errors by PBU-clustered bootstrap (500
replications) at fixed $\psi^*$.

Corner-solution tenders---convite auctions where the 3-bidder
minimum forces $m \geq \underline{n}(\tau) - n$---are excluded from
estimation and reserved for an overidentification check: the
interior-solution model should predict lower FL counts than observed
in these tenders, confirming that the constraint binds.

\paragraph{Identification.}
$\gamma$ is identified from the covariation of $\bar{m}$ with
$\bar{n}$ across PBUs (holding PBU size fixed) and is the
best-determined parameter ($R^2 = 0.585$ from the log-linear first
stage). Its sign is positive, consistent with the strategic
complementarity predicted by Proposition~\ref{prop:market_selection}:
cartels deploy more cover bidders in tenders with more genuine
competition. This relationship holds within PBU (PBU fixed effects:
$+0.29$) as well as between PBUs ($+0.45$), ruling out pure
cross-sectional confounding.
The level of $c_1$ is pinned down by the unconditional mean of $m^*$.
The ratio $\phi_0 / c_1$ is identified from the cross-PBU gradient of
$\bar{m}$ with PBU size (holding $\bar{n}$ fixed): larger PBUs with
higher detection probability should attract fewer cover bidders. In
practice, this gradient is economically small ($-0.018$, $p = 0.024$),
and we cannot reject that enforcement has zero effect on cover-bidder
deployment. This is itself a substantive finding: cartel cover-bidding
intensity responds to competitive pressure, not to oversight.

% ── 5.3 Reduced-Form Tests ──────────────────────────────────────

\subsection{Reduced-Form Tests of Structural Predictions}
\label{sec:reduced_form}

\paragraph{Reduced-form tests (Predictions~\ref{pred:price}--\ref{pred:selection}).}
\begin{equation}
  y_{igt} = \beta \cdot \text{losers}_{igt} + \mathbf{x}_{igt}'
  \boldsymbol{\delta} + \alpha_g + \lambda_t + \gamma_k + \varepsilon_{igt}
  \label{eq:ols_main}
\end{equation}
Standard errors clustered at the item level. Four specifications:
(1)~item + year FE; (2)~+ PBU FE; (3)~preg\~ao only; (4)~convite
only. Cross-fitting (FL defined on odd years, estimated on even, and
vice versa) breaks the classification--outcome link and yields a
conservative lower bound. CEM and IPW matching provide non-parametric
checks.

\paragraph{IV as measurement-error diagnostic.}
A leave-one-out instrument
$Z_{kgt} = \sum_{j \neq k} \mathbf{1}[\text{FL active at PBU } j,
\text{ group } g, \text{ year } t]$
exploits supply-side variation ($F = 396$). The IV is not our
primary estimator---balance tests reveal systematic imbalance across
observables (Appendix~\ref{app:robustness})---but it serves as a
measurement-error diagnostic: the point estimate (0.194) exceeds
OLS (0.064) by a factor of three, consistent with attenuation from
binary misclassification of FL status.

% ── 5.4 Additional Tests ──────────────────────────────────────────

\subsection{Specification and Detection Tests}
\label{sec:additional_tests}

\paragraph{Detection comparison.}
\label{sec:detection_validation}
We benchmark the FL screen against three alternatives using CADE
convictions as ground truth (temporal split:
Section~\ref{sec:cade_sample}): (i)~\citet{imhof2019detecting}-style
bid-level features; (ii)~a random forest (500 trees, 5-fold CV);
(iii)~the structural likelihood-ratio score.

\paragraph{Bajari--Ye tests.}
\label{sec:bajari_ye}
First-stage residuals from
$\log b_{jt} = \mathbf{x}_j'\boldsymbol{\phi} + \alpha_i +
\lambda_t + u_{jt}$, excluding $n_{\text{bids}}$
\citep[jointly determined;][p.~978]{bajari2003deciding}.
Tests: exchangeability (KS on FL vs.\ non-FL residuals) and
conditional independence (pairwise product of FL residuals). A
tender-FE variant tests the novel Regime~2 prediction: the FL
pairwise product should drop \emph{below} the non-FL product, as
the tender mean absorbs the shared focal-point component.

\paragraph{Market-selection heterogeneity.}
\label{sec:empirical_network}
For each FL firm, we compute the winner HHI
$= \sum_{w} s_{wj}^2$ and split at the median. Under
Proposition~\ref{prop:market_selection}, the FL--price association
should concentrate in competitive markets.
